% Part: methods
% Chapter: proofs
% Section: inference-patterns

\documentclass[../../../include/open-logic-section]{subfiles}

\begin{document}

\olfileid[tr]{mth}{prf}{pat}

\olsection{Çıkarım Kalıpları}

Kanıtlar tek tek çıkarımlardan oluşur. Bir çıkarım yaptığımızı çoğu
zaman ``öyleyse'', ``böylece'' ya da ``dolayısıyla'' gibi bir sözcükle
belirtiriz. Çıkarım genellikle kanıtımızda hâlihazırda bulunan bir ya da
iki olguya dayanır---bu, varsaydığımız bir şey ya da daha önce bir
çıkarımla ulaştığımız bir sonuç olabilir. Açık olmak için bu şeyleri
etiketleyebilir, çıkarım sırasında başka hangi ifadeleri kullandığımızı
belirtebiliriz. Bir çıkarım çoğu zaman, yeni sonucumuzun önceki
şeylerden \emph{neden} çıktığını açıklayan bir gerekçe de içerir.
Kanıtlarda çok sık kullanılan bazı yaygın çıkarım kalıpları vardır;
aşağıda bunlardan bazılarını ele alacağız. Tümevarımla kanıt gibi kimi
çıkarım kalıpları daha kapsamlıdır (ve daha sonra tartışılacaktır).

Bir çıkarım kalıbını daha önce ele aldık: bir tanımı açmak ya da
uygulamak. Bir tanımı açtığımızda, tanımlananı içeren bir şeyi
tanımlayanı kullanarak yeniden ifade ederiz. Örneğin bir kanıtın
akışı içinde $D = E$ (a) olduğunu zaten belirlediğimizi varsayalım.
O zaman kümeler için $=$ tanımını uygulayıp şu çıkarımı yapabiliriz:
``Öyleyse (a)'dan tanım gereği, $D$'de bulunan her !!{element}, $E$'de
de bir !!a{element} olarak bulunur ve bunun tersi de geçerlidir.''

Biraz kafa karıştırıcı biçimde, bir çıkarımın gerekçesini çoğu kez
çıkarımı gerçekten yaptığımız sırada değil, ondan önce yazarız. Henüz
$D = E$'yi kanıtlamadığımızı fakat kanıtlamak istediğimizi düşünün.
Hedeflediğimiz sonuç $D = E$ ise tanımı uygulayarak bu hedefi de yeniden
ifade edebiliriz: $D = E$'yi kanıtlamak için $D$'de bulunan her
!!{element} için $E$'de bir !!a{element} bulunduğunu ve bunun tersini
kanıtlamalıyız. Dolayısıyla kanıtımız şu biçimde olur: (a) $D$'deki her
!!{element} için $E$'de bir !!a{element} bulunduğunu kanıtla; (b)
$E$'deki her !!{element} için $D$'de bir !!a{element} bulunduğunu kanıtla; (c)
öyleyse (a) ve (b)'den, $=$ tanımı gereği, $D = E$. Ancak bunu genellikle
bu biçimde yazmayız. Bunun yerine şöyle bir şey yazabiliriz:
\begin{quote}
$D = E$ olduğunu göstermek istiyoruz. $=$ tanımı gereği bu, $D$'deki
her !!{element} için $E$'de bir !!a{element} bulunduğunu ve bunun tersini
göstermek demektir.

(a) \dots ($D$'deki her !!{element} için $E$'de bir !!a{element}
olduğunun kanıtı) \dots

(b) \dots ($E$'deki her !!{element} için
$D$'de bir !!a{element} bulunduğunun kanıtı) \dots
\end{quote}

\subsection{Bir Bağlaşımı Kullanmak}

Belki de en basit çıkarım kalıbı, bir bağlaşımın bileşenlerinden birini
sonuç olarak çıkarmaktır. Başka bir deyişle: $p$ ve~$q$ olduğunu
varsaymış ya da zaten kanıtlamışsak $p$ sonucunu (ve aynı zamanda $q$
sonucunu) çıkarabiliriz. Bu öylesine temel bir çıkarımdır ki çoğu kez
adı bile anılmaz. Örneğin $D = E$ tanımını açtığımızda $D$'de bulunan
her !!{element} için $E$'de bir !!a{element} bulunduğunu ve bunun
tersini belirlemiş oluruz. Buradan $E$'deki her !!{element} için $D$'de
bir !!a{element} bulunduğu sonucuna varabiliriz (``bunun tersi'' kısmı budur).

\subsection{Bir Bağlaşımı Kanıtlamak}

Bazen kanıtlamanız istenen şey bir bağlaşım biçiminde olur; sizden
``$p$ ve $q$'yu kanıtlamanız'' istenir. Bu durumda yalnızca iki şey
yapmalısınız: $p$'yi, ardından $q$'yu kanıtlamak. Kanıtınızı iki kısma
ayırıp açıklık için bunları etiketleyebilirsiniz. İlk notlarınızı
alırken sayfanın başına ``(1) $p$'yi kanıtla'', ortasına da ``(2)
$q$'yu kanıtla'' yazabilirsiniz. (Elbette sizden açıkça bir bağlaşımı
kanıtlamanız istenmemiş olabilir; yine de kanıtınızın bir bağlaşımı
kanıtlamanızı gerektirdiğini görebilirsiniz. Örneğin $D = E$'yi
kanıtlamanız istenirse $=$ tanımını açtıktan sonra şunları kanıtlamanız
gerektiğini görürsünüz: $D$'deki her !!{element}, $E$'de bir
!!a{element} olarak bulunur \emph{ve} $E$'deki her !!{element}, $D$'de
bir !!a{element} olarak bulunur.)

\subsection{Bir Ayrışımı Kanıtlamak}

Kanıtladığınız şey bir ayrışım biçimindeyse (yani ``$p$ ya da $q$''
biçiminde bir ifadeyse), ayrışanlardan birinin doğru olduğunu göstermek
yeterlidir. Ancak ayrışanlardan herhangi birinin teoreminizin
varsayımlarından doğrudan çıkması hemen hemen hiç görülmez. Daha sık
olarak, teoreminizin varsayımları da ayrışıktır ya da belirli türdeki
bütün şeylerin iki özellikten birine sahip olduğunu gösteriyorsunuzdur;
ama bazıları bir, bazıları öteki özelliğe sahiptir. Durumlara göre kanıt
burada işe yarar (aşağıya bakın).

\subsection{Koşullu Kanıt}

Karşılaşacağınız birçok teorem koşullu biçimdedir (yani $p$ geçerliyse
$q$'nun da doğru olduğunu göstermeniz istenir). Bu durumların kuruluşu
kolaydır---koşullunun önbileşenini (burada $p$'yi) varsayın ve buradan
$q$ sonucunu kanıtlayın. Dolayısıyla teoreminiz ``$p$ ise $q$'' diyorsa,
kanıtınıza ``$p$'yi varsayalım'' diye başlar ve sonunda~$q$'yu
kanıtlamış olursunuz.

Koşullular farklı biçimlerde ifade edilebilir. Bir teorem ``$p$ ise
$q$'' yerine ``$p$, ancak $q$ ise'', ``$q$, $p$ ise'' ya da ``$q$,
$p$ koşuluyla'' diyebilir. Bunların hepsi aynı anlama gelir ve $p$'yi
varsayıp bu varsayımdan~$q$'yu kanıtlamayı gerektirir. Bir iki yönlü
koşullunun (``$p$ ancak ve ancak $q$ ise'') aslında bir araya getirilmiş
iki koşullu olduğunu hatırlayın: $p$ ise $q$ ve $q$ ise~$p$. O hâlde
yapmanız gereken yalnızca iki koşullu kanıttır: biri ilk koşullu, diğeri
ikinci koşullu için. Bununla birlikte bazen bir ``ancak ve ancak''
ifadesini, ``$p$'' ile başlayıp ``$q$'' ile bitecek biçimde başka birkaç
``ancak ve ancak'' ifadesini zincirleyerek kanıtlamak mümkündür---ama bu
durumda her adımın gerçekten bir ``ancak ve ancak'' olduğundan emin
olmalısınız.

\subsection{Tümel İddialar}

Tümel bir iddiayı kullanmak basittir: bir şey her nesne için doğruysa,
tek tek her nesne için de doğrudur. Örneğin kanıtınızın hipotezi $A
\subseteq B$ ise bu, $\subseteq$ tanımı açıldığında, her $x \in A$ için
$x \in B$ olduğu anlamına gelir. Dolayısıyla $z \in A$ olduğunu zaten
biliyorsanız $z \in B$ sonucunu çıkarabilirsiniz.

Tümel bir iddiayı kanıtlamak biraz çetrefilli görünebilir. Genellikle
bu ifadeler şu biçimdedir: ``$x$, $P$ özelliğine sahipse $Q$ özelliğine
de sahiptir'' ya da ``Bütün $P$'ler $Q$'dur.'' Elbette ifade bu biçime
tam olarak uymayabilir ve sizden tam olarak neyin kanıtlanmasının
istendiğini anlamak biraz alıştırma gerektirir. Ama çoğu zaman, belli
bir özelliğe sahip bütün nesnelerin başka bir belirli özelliğe sahip
olduğunu kanıtlamamız gerekir.

Tümel bir iddiayı kanıtlamanın yolu, ilk özelliğe sahip şeyler için
adlar ya da değişkenler ortaya koymak ve sonra onların öteki özelliğe
de sahip olduğunu göstermektir. Bunu, \emph{bütün}~$P$'ler için bir şeyi
kanıtlamak üzere onu \emph{keyfî} bir~$P$ için kanıtlamanız gerekir,
diye ifade edebiliriz. Ortaya konan ad da keyfî bir~$P$'nin adıdır.
Keyfî şeyler için ad olarak genellikle tek harf kullanır ve çoğu zaman
yerleşik kurallara uyarız: örneğin doğal sayılar için $n$,
!!{formula}s için $!A$, kümeler için $A$, fonksiyonlar için $f$ vb.

İncelik, kanıt boyunca genelliği korumaktır. Keyfî bir nesnenin
(``$x$'') $P$ özelliğine sahip olduğunu varsayarak başlar ve (yalnızca
tanımlara ya da varsaymanıza izin verilenlere dayanarak) $x$'in $Q$
özelliğine sahip olduğunu gösterirsiniz. $x$'in, $P$ özelliğine sahip
olması dışında özellikle ne olduğunu belirtmediğiniz için, $P$'ye
sahip her şeyin $Q$ özelliğine sahip olduğunu ileri sürebilirsiniz.
Kısacası $x$, $P$ özelliğine sahip \emph{bütün} şeylerin yerini tutar.

\begin{prop}
  Bütün $A$ ve $B$ kümeleri için $A \subseteq A \cup B$.
\end{prop}

\begin{proof}
  $A$ ve $B$ keyfî kümeler olsun. $A \subseteq A \cup B$ olduğunu
  göstermek istiyoruz. $\subseteq$ tanımı gereği bu şu demektir: her
  $x$ için, $x \in A$ ise $x \in A \cup B$. Öyleyse $x \in A$, $A$'nın
  keyfî bir !!{element} olsun. $x \in A \cup B$
  olduğunu göstermeliyiz.
  $x \in A$ olduğundan, $x \in A$ ya da $x \in B$. Dolayısıyla $x \in
  \Setabs{x}{x \in A \lor x \in B}$. Ancak bu, $\cup$ tanımı gereği
  $x \in A \cup B$ demektir.
\end{proof}


\subsection{Durumlara Göre Kanıt}

Bir varsayım ya da daha önce belirlenmiş bir sonuç olarak elinizde bir
ayrışım bulunduğunu düşünün---$p$ ya da $q$'nun doğru olduğunu varsaymış
veya kanıtlamışsınız. $r$'yi kanıtlamak istiyorsunuz. Bunu iki adımda
yaparsınız: önce $p$'nin doğru olduğunu varsayıp~$r$'yi kanıtlar, sonra
$q$'nun doğru olduğunu varsayıp~$r$'yi yeniden kanıtlarsınız. Bu yöntem
işler; çünkü iki seçenekten birinin geçerli olduğunu varsayıyor ya da
biliyoruzdur. İki adım da seçeneklerden her birinin~$r$'nin doğruluğu
  için yeterli olduğunu belirler. (İkisi de doğruysa $r$'nin doğru olması
  için bir değil iki gerekçemiz vardır. $r$'nin hem $p$ hem de~$q$
  varsayımı altında doğru olduğunu ayrıca kanıtlamak gerekmez.) Ne
yaptığımızı belirtmek için ``durumları ayırdığımızı'' duyururuz.
Örneğin $x \in B \cup C$ olduğunu bildiğimizi varsayalım. $B \cup C$,
$\Setabs{x}{x \in B \text{ ya da } x \in C}$ olarak tanımlanır. Başka
bir deyişle, tanım gereği $x \in B$ ya da $x \in C$. Buradan $x \in A$
olduğunu kanıtlamak için önce $x \in B$'yi varsayıp bu varsayımdan $x
\in A$'yı kanıtlar, sonra $x \in C$'yi varsayıp $x \in A$'yı yeniden
kanıtlardınız. Varsayımın altına ``Durumları ayıralım'', onun altına
``Durum (1): $x \in B$'', sayfanın ortasına da ``Durum (2): $x \in C$''
yazardınız. Sonra sayfanın üst ve alt yarılarını doldurmaya geçerdiniz.

Durumlara göre kanıt, kanıtladığınız şeyin kendisi ayrışık olduğunda
özellikle kullanışlıdır. İşte basit bir örnek:

\begin{prop}
$B \subseteq D$ ve $C \subseteq E$ olduğunu varsayın. O zaman $B \cup C
\subseteq D \cup E$.
\end{prop}

\begin{proof}
  (a) $B \subseteq D$ ve (b) $C \subseteq E$ olduğunu varsayalım. Tanım
  gereği her $x \in B$, aynı zamanda $\in D$'dir (c) ve her $x \in C$,
  aynı zamanda $\in E$'dir (d). $B \cup C \subseteq D \cup E$ olduğunu
  göstermek için, $x \in B \cup C$ ise $x \in D \cup E$ olduğunu
  göstermeliyiz ($\subseteq$ tanımı gereği). $x \in B \cup C$, ancak ve
  ancak $x \in B$ ya da $x \in C$ ise geçerlidir ($\cup$ tanımı gereği).
  Benzer biçimde, $x \in D \cup E$, ancak ve ancak $x \in D$ ya da $x
  \in E$ ise geçerlidir. Öyleyse şunu göstermeliyiz: herhangi bir $x$
  için, $x \in B$ ya da $x \in C$ ise $x \in D$ ya da $x \in E$.

  \begin{quote}
  Şimdiye kadar yalnızca tanımları açtık!{} Önermemizi $\subseteq$ ve
  $\cup$ kullanmadan yeniden ifade ettik ve geriye tümel bir koşullu
  iddiayı kanıtlamak kaldı. Yukarıdaki tartışmamıza göre bunu, ``ise''
  kısmının doğru olduğunu varsaydığımız bir $x$ alıp ``o hâlde''
  kısmının da doğru olduğunu göstererek yaparız. Başka bir deyişle $x
  \in B$ ya da $x \in C$ olduğunu varsayıp $x \in D$ ya da $x \in E$
  olduğunu göstereceğiz.\footnote{Bu paragraf yalnızca ne yaptığımızı
    açıklıyor---kanıtın bir parçası değildir ve kendi kanıtlarınızı
    yazarken bu kadar ayrıntıya girmeniz gerekmez.}
  \end{quote}

  $x \in B$ ya da $x \in C$ olduğunu varsayalım. $x \in D$ ya da $x \in
  E$ olduğunu göstermeliyiz. Durumları ayıralım.

  Durum 1: $x \in B$. (c)'ye göre $x \in D$. Öyleyse $x \in D$ ya da
  $x \in E$. (Burada bir önceki alt kısımda ele alınan çıkarımı yaptık!)

  Durum 2: $x \in C$. (d)'ye göre $x \in E$. Öyleyse $x \in D$ ya da
  $x \in E$.
 \end{proof}


\subsection{Bir Varlık İddiasını Kanıtlamak}

Bir varlık iddiasını kanıtlamanız istendiğinde soru genellikle
``bir~$x$ için $\dots x \dots$ olduğunu kanıtlayın'', yani
``$\dots x \dots$'' ile betimlenen özelliğe sahip bir nesnenin var
olduğunu kanıtlayın biçiminde olur. Bu durumda uygun bir nesne belirleyip
onun gereken özelliğe sahip olduğunu göstermelisiniz. Bu kolay
görünebilir, fakat bu tür bir kanıt çetrefilli olabilir. Genellikle bir
nesneyi \emph{kurmayı} ya da \emph{tanımlamayı} ve böyle tanımlanan
nesnenin gereken özelliğe sahip olduğunu kanıtlamayı içerir. Doğru
nesneyi bulmak zor olabilir, gereken özelliğe sahip olduğunu kanıtlamak
zor olabilir, hatta bazen herhangi bir nesneyi tanımlamayı başardığınızı
göstermek bile çetrefillidir!{}

Genellikle bunu nesneyi belirterek, örneğin ``$x$, \dots{} olsun''
(burada \dots{} aklınızdaki nesneyi belirler) diye yazar; belki
$\dots$'ın gerçekten var olan bir nesneyi betimlediğini kanıtlar ve sonra
$x$'in $Q$ özelliğine sahip olduğunu göstermeye geçersiniz. İşte basit
bir örnek.

\begin{prop}
  $x \in B$ olduğunu varsayın. O zaman bir~$A$ vardır ve $A \subseteq B$
  ile $A \neq \emptyset$ geçerlidir.
\end{prop}

\begin{proof}
  $x \in B$ olduğunu varsayalım. $A = \{x\}$ olsun.
  \begin{quote}
    Burada~$A$ kümesini, !!{element}s sıralamasını vererek tanımladık. $x$'in bir
    nesne olduğunu varsaydığımızdan ve !!{element}s sıralamasını vererek her
    zaman bir küme oluşturabildiğimizden, burada bir~$A$ kümesi
    tanımlamayı başardığımızı göstermemiz gerekmez. Bununla birlikte
    $A$'nın önermenin istediği özelliklere sahip olduğunu hâlâ
    göstermeliyiz. Kanıt onsuz tamamlanmış değildir!{}
  \end{quote}
  $x \in A$ olduğundan $A \neq \emptyset$.
  \begin{quote}
    Bu, $A$'nın $\{x\}$ olarak tanımlanmasına ve $x \in \{x\}$ ile $x
    \notin \emptyset$ şeklindeki açık olgulara dayanır.
  \end{quote}
  $x$, ~$\{x\}$ içinde bulunan tek nesnedir; dolayısıyla bu kümenin tek
  !!{element} odur. Ayrıca $x \in B$'dir; böylece $A$'daki her
  !!{element}, $B$'de bir !!a{element} olarak
  bulunur. $\subseteq$
  tanımı gereği $A \subseteq B$.
\end{proof}

\subsection{Varlık İddialarını Kullanmak}

Bir varlık iddiasının doğru olduğunu bildiğinizi (onu kanıtladığınızı
ya da kullanabileceğiniz bir hipotez olduğunu) düşünün; örneğin
``bir~$x$ için $x \in A$'' ya da ``bir $x \in A$ vardır.'' Bunu
kanıtınızda kullanmak istiyorsanız, hipotezinizin var olduğunu söylediği
şeylerden biri için bir adınız varmış gibi davranabilirsiniz. $A$ en az
bir şey içerdiğine göre, bu adın gönderimde bulunabileceği şeyler
vardır. Elbette bunlardan birini seçemeyebilir ya da daha ayrıntılı
betimleyemeyebilirsiniz ($\in A$ olması dışında). Fakat kanıtın amacı
bakımından onu seçmiş gibi davranıp ona bir ad verebilirsiniz. Daha önce
kullanmadığınız (ya da hipotezlerinizde geçmeyen) bir ad seçmek
önemlidir; aksi hâlde işler ters gidebilir. Kanıtınızda bunu ``bir $x$
için $x \in A$''dan ``$a \in A$ olsun''a geçerek belirtirsiniz. Artık
~$a$ hakkında akıl yürütebilir, başka hipotezler kullanabilir ve bir $p$
sonucuna ulaşıncaya dek devam edebilirsiniz. $p$ artık~$a$'dan söz
etmiyorsa, $p$, $a \in A$ varsayımından bağımsızdır ve yalnızca ``bir
$x$ için $x \in A$'' varsayımından çıktığını göstermişsinizdir.

\begin{prop}
$A \neq \emptyset$ ise $A \cup B \neq \emptyset$.
\end{prop}

\begin{proof}
  $A \neq \emptyset$ olduğunu varsayalım. Öyleyse bir $x$ için $x \in A$.
  \begin{quote}
    Burada önce yalnızca önermenin hipotezini yeniden ifade ettik. Bu
    hipotez, yani $A \neq \emptyset$, yalnızca birkaç tanımı açarak
    ulaşabileceğiniz bir varlık iddiasını gizler. $=$ tanımı bize $A =
    \emptyset$'in ancak ve ancak her $x \in A$ aynı zamanda $\in
    \emptyset$ ve her $x \in \emptyset$ aynı zamanda $\in A$ ise geçerli
    olduğunu söyler. İki tarafı da olumsuzlayınca şunu elde ederiz: $A
    \neq \emptyset$, ancak ve ancak bir $x \in A$, $\notin \emptyset$
    ise ya da bir $x \in \emptyset$, $\notin A$ ise geçerlidir. Hiçbir
    şey $\in \emptyset$ olmadığından ikinci ayrışan asla doğru olamaz;
    ``$x \in A$ ve $x \notin \emptyset$'' de yalnızca $x \in A$'ya
    indirgenir. Dolayısıyla $A \neq \emptyset$, ancak ve ancak bir $x$
    için $x \in A$ ise geçerlidir. Bu bir varlık iddiasıdır. Şimdi bu
    varlık iddiasını, $A$'da bulunan !!{element}s arasından birine ad vererek kullanırız:
  \end{quote}
  $a \in A$ olsun.
  \begin{quote}
    Şimdi~$\in A$ olan şeylerden birine bir ad verdik. $a$ hakkında
    akıl yürütmeye devam edeceğiz ama $a \in A$ dışında hiçbir şey
    varsaymamaya dikkat edeceğiz:
  \end{quote}
  $a \in A$ olduğundan $a \in A \cup B$; bu, $\cup$ tanımı gereğidir.
  Dolayısıyla bir $x$ için $x \in A \cup B$, yani $A \cup B \neq
  \emptyset$.
  \begin{quote}
    Son adımda ``$a \in A \cup B$''den ``bir $x$ için $x \in A \cup
    B$''ye geçtik. İkinci ifade artık $a$'dan söz etmez; böylece ``bir
    $x$ için $x \in A \cup B$''nin yalnızca ``bir $x$ için $x \in A$''dan
    çıktığını biliriz. Bu da $A \cup B \neq \emptyset$ demektir.
  \end{quote}
\end{proof}

Bağlı değişkenleri, yaptığımız gibi, ``$x$'' gibi yazıp varsayımsal
adlar olan $a$ gibi adlardan ayrı tutmak iyi bir uygulama olabilir.
Ancak pratikte çoğu zaman bunu yapmaz ve yalnızca $x$ kullanırız:
\begin{quote}
$A \neq \emptyset$ olduğunu, yani bir $x \in A$ bulunduğunu varsayalım.
$\cup$ tanımı gereği $x \in A \cup B$. Öyleyse $A \cup B \neq
\emptyset$.
\end{quote}
Ancak bunu yaptığınızda, farklı varlık iddiaları için farklı $x$ ve
$y$'ler kullanmaya özellikle dikkat etmelisiniz. Örneğin aşağıdaki
\emph{doğru bir kanıt değildir}: ``$A \neq \emptyset$ ve $B \neq
\emptyset$ ise $A \cap B \neq \emptyset$'' (bu ifade zaten doğru değildir).
\begin{quote}
$A \neq \emptyset$ ve $B \neq \emptyset$ olduğunu varsayalım. Öyleyse
bir $x$ için $x \in A$ ve ayrıca bir $x$ için $x \in B$. $x \in A$ ve
$x \in B$ olduğundan $x \in A \cap B$; bu, $\cap$ tanımı gereğidir.
Dolayısıyla $A \cap B \neq \emptyset$.
\end{quote}
Yanlış adımın nerede olduğunu bulup sonucun neden geçerli olmadığını
açıklayabilir misiniz?

\end{document}
